\section*{Hoofdstuk \ref{ch:g10:numbers} --- Getallen en getallenverzamelingen}

\begin{solution}{exo:g10:numbers:1}
$\frac{15}{3} = 5 \in \N$. \quad $-7 \in \Z$. \quad
$\frac{22}{7} \in \Q$ (het is geen \omterm{def:g10:numbers:sets}{geheel getal}: $22 = 7 \times 3 + 1$).
\quad $\sqrt 9 = 3 \in \N$. \quad $\sqrt{10} \in \R$ (\omterm{ex:g10:numbers:classify}{irrationaal}, want
$10$ is het kwadraat van geen enkel \omterm{def:g10:numbers:sets}{rationaal getal} --- hier aangenomen, in
de geest van \cref{thm:g10:numbers:sqrt2}). \quad
$-2.4 = -\frac{24}{10} \in \Q$. \quad $0 \in \N$.
\end{solution}

\begin{solution}{exo:g10:numbers:2}
(a) $\intco{-2}{5}$: volle stip in $-2$, open stip in $5$.
(b) $\intoo{3}{+\infty}$: open stip in $3$, arcering naar rechts.
(c) $\intoc{-\infty}{-1}$: arcering van links tot een volle stip in $-1$.
(d) ``afstand van $x$ tot $2$ hoogstens $3$'' betekent $\abs{x - 2} \leq 3$,
dus $x \in \intcc{-1}{5}$: volle stippen in $-1$ en $5$.
\end{solution}

\begin{solution}{exo:g10:numbers:3}
(a) De twee intervallen overlappen tussen $0$ en $2$:
$I \cap J = \intcc{0}{2}$ ($0 \in J$ en $2 \in I$, en beide behoren ook tot
de andere verzameling), en $I \cup J = \intco{-3}{4}$.

(b) $I \cap J$ is de verzameling van de $x$ met $-2 \leq x$ en $x < 1$:
$\intco{-2}{1}$. De \omterm{def:g10:numbers:interunion}{vereniging} bedekt alles: $I \cup J = \R$.
\end{solution}

\begin{solution}{exo:g10:numbers:4}
$\abs{-6} = 6$; \quad $\abs{4 - 9} = \abs{-5} = 5$; \quad
$\abs{-3 - 5} = \abs{-8} = 8$; \quad
$\sqrt 2 > 1$, dus $\abs{\sqrt2 - 1} = \sqrt2 - 1$; \quad
$1 - \sqrt2 < 0$, dus ook $\abs{1 - \sqrt2} = \sqrt2 - 1$ (een getal en zijn
tegengestelde hebben dezelfde \omterm{def:g10:numbers:abs}{absolute waarde}).
\end{solution}

\begin{solution}{exo:g10:numbers:5}
$\abs{x} = 5$: afstand $5$ tot $0$, dus $x = 5$ of $x = -5$.

$\abs{x - 1} = 3$: afstand $3$ tot $1$, dus $x = 4$ of $x = -2$.

$\abs{x - 4} \leq 1$: afstand hoogstens $1$ tot $4$, dus
$x \in \intcc{3}{5}$.

$\abs{x + 2} < 3$: afstand strikt kleiner dan $3$ tot $-2$, dus
$x \in \intoo{-5}{1}$.
\end{solution}

\begin{solution}{exo:g10:numbers:6}
Elk \omterm{def:g10:numbers:interval}{interval} wordt beschreven door zijn middelpunt $a$ (het midden van de
grenzen) en zijn straal $r$ (de halve lengte).

$\intcc{2}{8}$: $a = 5$, $r = 3$, dus $\abs{x - 5} \leq 3$.

$\intoo{-5}{1}$: $a = -2$, $r = 3$, open grenzen, dus $\abs{x + 2} < 3$.

$\intcc{-7}{-3}$: $a = -5$, $r = 2$, dus $\abs{x + 5} \leq 2$.
\end{solution}

\begin{solution}{exo:g10:numbers:7}
Zij $x = 0.272727\dots$ Dan is $100x = 27.2727\dots$, en aftrekken geeft
\[
100x - x = 27.2727\dots - 0.2727\dots = 27,
\]
dus $99x = 27$ en $x = \frac{27}{99} = \frac{3}{11}$, een quotiënt van
\omterm{def:g10:numbers:sets}{gehele getallen}: $x$ is rationaal.
\end{solution}

\begin{solution}{exo:g10:numbers:8}
\emph{1. Waar:} \omterm{def:g10:numbers:sets}{gehele getallen} (positieve of negatieve hele getallen)
optellen levert altijd een \omterm{def:g10:numbers:sets}{geheel getal}.

\emph{2. Niet waar:} $\frac{1}{2}$ is het quotiënt van de \omterm{def:g10:numbers:sets}{gehele getallen}
$1$ en $2$ en is zelf niet geheel.

\emph{3. Waar:} $\frac pq + \frac{p'}{q'} = \frac{pq' + p'q}{qq'}$ is
opnieuw een quotiënt van \omterm{def:g10:numbers:sets}{gehele getallen} (met noemer verschillend van nul).

\emph{4. Waar:} stel dat $r$ rationaal is, $t$ \omterm{ex:g10:numbers:classify}{irrationaal}, en dat
$r + t = s$ rationaal zou zijn. Dan zou $t = s - r$ een verschil van twee
\omterm{def:g10:numbers:sets}{rationale getallen} zijn, en dus rationaal (volgens 3, toegepast met $-r$)
--- tegenspraak. Bijgevolg is $r + t$ \omterm{ex:g10:numbers:classify}{irrationaal}.
\end{solution}

\begin{solution}{exo:g10:numbers:9}
Vermenigvuldigen van $1.414 \leq \sqrt2 \leq 1.415$ met $2 > 0$:
$2.828 \leq 2\sqrt2 \leq 2.830$.

$3$ optellen: $4.414 \leq \sqrt2 + 3 \leq 4.415$.

Vermenigvuldigen met $-1 < 0$ keert de ongelijkheden om:
$-1.415 \leq -\sqrt2 \leq -1.414$.
\end{solution}

\begin{solution}{exo:g10:numbers:10}
Eerst het hulpfeit. Deel $p$ door $3$: de rest is $0$, $1$ of $2$, dus
$p = 3k$, $p = 3k+1$ of $p = 3k+2$. Kwadrateren geeft
\[
(3k)^2 = 3(3k^2), \quad
(3k+1)^2 = 3(3k^2 + 2k) + 1, \quad
(3k+2)^2 = 3(3k^2 + 4k + 1) + 1 .
\]
Alleen het eerste is een veelvoud van $3$: is $p^2$ een veelvoud van $3$,
dan is $p$ dat ook.

Stel nu dat $\sqrt3 = \frac pq$ volledig vereenvoudigd. Kwadrateren geeft
$p^2 = 3q^2$, dus is $p^2$ een veelvoud van $3$, dus $p = 3k$. Dan is
$9k^2 = 3q^2$, dus $q^2 = 3k^2$ en is ook $q$ een veelvoud van $3$ --- maar
dan stond de breuk $\frac pq$ niet volledig vereenvoudigd: tegenspraak.
Bijgevolg is $\sqrt3$ \omterm{ex:g10:numbers:classify}{irrationaal}.
\end{solution}

\begin{solution}{pb:g10:numbers:1}
\textbf{1.} $-7 \in \Z$; \ $\frac{13}{4} \in \Q$; \
$\sqrt{16} = 4 \in \N$; \ $0.121212\ldots = \frac{12}{99} = \frac{4}{33}
\in \Q$ (de weekendopgave over repeterende decimalen in het
onderbouwvolume); \ $\sqrt8 = 2\sqrt2$ is \omterm{ex:g10:numbers:classify}{irrationaal}: kleinste
verzameling $\R$; \ $\pi$: $\R$.

\textbf{2.} $\frac pq + \frac rs = \frac{ps + qr}{qs}$ en
$\frac pq \times \frac rs = \frac{pr}{qs}$: quotiënten van \omterm{def:g10:numbers:sets}{gehele getallen}
met noemer verschillend van nul --- rationaal.

\textbf{3.} (a) Is $r$ rationaal, $x$ \omterm{ex:g10:numbers:classify}{irrationaal}, en zou $r + x = q$
rationaal zijn, dan was $x = q - r$ een verschil van \omterm{def:g10:numbers:sets}{rationale getallen}, en
dus rationaal (vraag 2): tegenspraak. (b) Is $r \neq 0$ en $r x = q$
rationaal, dan is $x = \frac qr$ rationaal: tegenspraak.

\textbf{4.} $\sqrt2$ en $-\sqrt2$ zijn \omterm{ex:g10:numbers:classify}{irrationaal} met som $0$; $\sqrt2$ en
$\sqrt2$ hebben product $2$. Rationale uitkomsten uit irrationale
ingrediënten: de irrationale getallen vormen geen stabiel koninkrijk.

\textbf{5.} $\sqrt8 = 2\sqrt2$, dus $\sqrt2 + \sqrt8 = 3\sqrt2$:
\omterm{ex:g10:numbers:classify}{irrationaal} (vraag 3b, met factor $3$). En
$\sqrt2 \times \sqrt8 = \sqrt{16} = 4 \in \N$: het product van twee
irrationale getallen belandt in het allerkleinste koninkrijk.

\textbf{6.} $3.475$, daarna $3.471$ (of $3.4701$, $3.47001$, \dots): door
decimalen aan te plakken maak je eindeloos nieuwe \omterm{def:g10:numbers:sets}{rationale getallen} tussen
de twee.

\textbf{7.} De veelvouden van $10^{-n}$ marcheren over de lijn met stappen
van $10^{-n} < g$. Het eerste veelvoud strikt groter dan $a$ --- het
bestaat, want de veelvouden overtreffen $a$ uiteindelijk --- ligt hoogstens
één stap voorbij $a$, dus vóór $a + g = b$: strikt tussen $a$ en $b$. Een
roosterpunt is een decimaal getal en dus rationaal: elk \omterm{def:g10:numbers:interval}{interval} van
positieve lengte bevat er een.

\textbf{8.} Tussen $a$ en het rationale getal $r_1$ uit vraag 7 ligt
(opnieuw vraag 7) een \omterm{def:g10:numbers:sets}{rationaal getal} $r_2$; tussen $a$ en $r_2$ een
\omterm{def:g10:numbers:sets}{rationaal getal} $r_3$; enzovoort: oneindig veel, alle verschillend, alle in
het oorspronkelijke \omterm{def:g10:numbers:interval}{interval}.

\textbf{9.} $r + \frac{\sqrt2}{10^k}$ is de som van een rationaal en een
\omterm{ex:g10:numbers:classify}{irrationaal getal} ($\frac{\sqrt2}{10^k}$ is \omterm{ex:g10:numbers:classify}{irrationaal} volgens vraag 3b),
en dus \omterm{ex:g10:numbers:classify}{irrationaal}. Omdat $\frac{\sqrt2}{10^k} < \frac{2}{10^k}$ onder elke
grens zakt, ligt $r + \frac{\sqrt2}{10^k}$ voor $k$ groot genoeg nog steeds
vóór $b$: een \omterm{ex:g10:numbers:classify}{irrationaal getal} binnen het \omterm{def:g10:numbers:interval}{interval} --- en $k$ laten
variëren geeft er oneindig veel.

\textbf{10.} Er is geen kleinste positief \omterm{def:g10:numbers:sets}{reëel getal}: was $s > 0$ dat wel,
dan zou het \omterm{def:g10:numbers:interval}{interval} $(0, s)$ nog altijd een \omterm{def:g10:numbers:sets}{rationaal getal} bevatten
(vraag 7), positief en kleiner dan $s$. Er is ook geen getal vlak na $3$:
elke kandidaat $c > 3$ laat het \omterm{def:g10:numbers:interval}{interval} $(3, c)$ over, en dat is niet leeg
--- het spel uit de weekendopgave over het ontbreken van een volgend getal
in het onderbouwvolume, nu een stelling over $\R$.

\textbf{11.} $\abs{x - 5} = 2$: $x = 3$ of $x = 7$.
$\abs{x - 5} < 2$: $x \in \intoo{3}{7}$.
$\abs{x + 1} \geq 3$: afstand tot $-1$ minstens $3$, dus
$x \in \intoc{-\infty}{-4} \cup \intco{2}{+\infty}$.

\textbf{12.} $\abs{d - 80} \leq 0.05$, dus
$d \in \intcc{79.95}{80.05}$. De zuiger van $79.97$ wordt goedgekeurd; die
van $80.06$ valt af, met een honderdste millimeter.

\textbf{13.} $(3, -5)$: $\abs{-2} = 2 \leq 8$. $(-2, -7)$:
$\abs{-9} = 9 = 2 + 7$: gelijkheid. Zelfde teken: $\abs{a + b}$ is de som
van de afstanden, gelijk aan $\abs a + \abs b$. Tegengesteld teken: de
wandeling keert terug, $\abs{a + b}$ is het \emph{verschil} van de
afstanden, strikt kleiner dan hun som (tenzij een van beide $0$ is).
Gelijkheid precies dan als $a$ en $b$ hetzelfde teken hebben of een van
beide nul is.

\textbf{14.} De oplossingen zijn de punten op gelijke afstand van $2$ en
$-4$: het midden, $x = -1$. Het is de middelloodlijn uit de weekendopgave
over de twee spiegels in het onderbouwvolume, ineengeschoven tot één
dimensie: een enkel punt.

\textbf{15.} $\sqrt{x^2} = \abs{x}$, niet $x$: voor $x = -3$ is
$\sqrt{9} = 3 = \abs{-3} \neq -3$. Dan leest $x^2 < 9$ als
$\sqrt{x^2} < 3$, dus $\abs x < 3$: $x \in \intoo{-3}{3}$.

\textbf{16.} De meting beweert alleen dat de lengte in het \omterm{def:g10:numbers:interval}{interval}
$\intcc{1.41420}{1.41422}$ ligt --- en volgens de vragen 7 en 9 bevat
\emph{dat \omterm{def:g10:numbers:interval}{interval} oneindig veel rationale en oneindig veel irrationale
getallen}. Hetzelfde geldt voor elke toekomstige tolerantie, hoe klein ook.
Geen enkele meting kan de twee families ooit scheiden; alleen een bewijs
over de exacte lengte kan dat --- zoals dat voor de diagonaal van het
eenheidsvierkant (de weekendopgave over irrationaliteit in het
onderbouwvolume).

\textbf{17.} Ze naderen $\sqrt2$ met fouten onder $10^{-1}, 10^{-2},
10^{-3}, \dots$: \omterm{def:g10:numbers:sets}{rationale getallen} dringen op elke schaal tegen $\sqrt2$
aan. Het algemene feit: elk \omterm{def:g10:numbers:sets}{reëel getal} wordt zo dicht als je wilt benaderd
door \omterm{def:g10:numbers:sets}{rationale getallen} (zijn decimale afkappingen) --- opnieuw dichtheid,
nu gezien vanuit het doel.

\textbf{18.} $a + b \in \intcc{2.4 + 1.1}{2.6 + 1.3} = \intcc{3.5}{3.9}$:
onzekerheid $\pm 0.2$ (de fouten tellen op).
$a \times b \in \intcc{2.4 \times 1.1}{2.6 \times 1.3} =
\intcc{2.64}{3.38}$: een onzekerheid van ongeveer $\pm 0.37$ rond $3$ ---
bij producten tellen de relatieve fouten op, zodat de nauwkeurigheid
sneller verslechtert.

\textbf{19.} Schrijf de ware waarden als $a = 2.5 + e_1$ en
$b = 1.2 + e_2$ met $\abs{e_1}, \abs{e_2} \leq 0.1$. Dan is de fout op de
som $\abs{e_1 + e_2} \leq \abs{e_1} + \abs{e_2} \leq 0.2$: de vuistregel
van de ingenieur \emph{is} de driehoeksongelijkheid van vraag 13.

\textbf{20.} De reële lijn heeft geen gaten en geen buren: na een getal
komt geen ``volgend'' getal (vraag 10). Haar punten vallen uiteen in de
\omterm{def:g10:numbers:sets}{rationale getallen} --- precies de afbrekende of repeterende decimale
getallen (de weekendopgave daarover in het onderbouwvolume) --- en de
irrationale getallen, en de twee families zijn zo strak verweven dat elk
\omterm{def:g10:numbers:interval}{interval} er oneindig veel van elke soort bevat (vragen 8 en 9); daarom kan
geen meting, maar alleen een bewijs, ze uit elkaar houden (vraag 16). En de
laatste verrassing, als belofte achtergelaten: er zijn meer irrationale dan
\omterm{def:g10:numbers:sets}{rationale getallen} --- niet door één voor één te tellen, maar in de
precieze zin van het vergelijken van oneindigheden, een theorie die in de
universitaire volumes wordt opgebouwd.
\end{solution}
